%% ============================================================
%% LUNA: Scalable Evidence-Based Oracles Using Bayesian Networks
%% Target: IEEE Journal of Biomedical and Health Informatics
%% ============================================================
%% COMPILE: pdflatex -> bibtex -> pdflatex -> pdflatex
%%
%% REVIEW MARKERS (remove all before final submission):
%%   \placeholder{X} = blue  = fill with real experimental value
%%   \todo{X}        = red   = structural task remaining
%%   \addcite        = orange = missing citation to resolve
%% ============================================================

\documentclass[journal]{IEEEtran}

%% --- Packages ---
\usepackage[T1]{fontenc}
\usepackage[utf8]{inputenc}
\usepackage{amsmath,amssymb,amsthm}
\usepackage{booktabs}
\usepackage{graphicx}
\usepackage{xcolor}
\usepackage{hyperref}
\usepackage{url}
\usepackage{cite}
\usepackage{multirow}
\usepackage{array}
\usepackage{tabularx}
\usepackage{balance}

%% --- Review helpers (remove before submission) ---
\newcommand{\todo}[1]{\textcolor{red}{\textbf{[TODO: #1]}}}
\newcommand{\placeholder}[1]{\textcolor{blue}{\textbf{[#1]}}}
\newcommand{\addcite}{\textcolor{orange}{\textbf{[CITE]}}}

%% --- Theorem environments ---
\newtheorem{definition}{Definition}
\newtheorem{proposition}{Proposition}

%% ============================================================
\begin{document}

\title{LUNA: Scalable Evidence-Based Oracles Using\\
       Bayesian Networks for Healthcare Claim Verification}

%% TODO: Replace with real author names and affiliations.
%% For blind review, replace entire \author block with:
%% \author{Anonymous Authors --- Submitted for Review}
\author{First~Author,~\IEEEmembership{Member,~IEEE,}
        Second~Author,~\IEEEmembership{Senior Member,~IEEE,}
        and~Third~Author%
\thanks{Manuscript received \today.}%
\thanks{F. Author is with the Department of ...,
        University of ..., City, Country
        (e-mail: author@university.edu).}%
\thanks{S. Author and T. Author are with ...}}

\maketitle

%% ============================================================
%% ABSTRACT
%% Target: ~220 words. Five required elements:
%%   (1) Health informatics problem
%%   (2) Gap in existing work
%%   (3) Method and architecture
%%   (4) Concrete numerical results  <- fill \placeholder blocks
%%   (5) Clinical and domain significance
%% TODO: Fill ALL \placeholder{} after experiments are complete.
%% ============================================================
\begin{abstract}
Blockchain-based smart contracts are increasingly proposed
for healthcare claim verification, yet existing oracle designs
reduce heterogeneous clinical evidence to opaque scalar values,
obscuring how individual evidence sources contribute to
decisions and preventing meaningful audit.
Bayesian network~(BN)-based oracles address this by making
probabilistic reasoning explicit, but on-chain BN inference
incurs gas costs that grow with model size, limiting practical
deployability.
We present LUNA (Ledger-grounded Uncertainty-aware Networked
Architecture), a BN-based oracle architecture that resolves
this tension by anchoring BN parameters and evidence
assignments on-chain while executing inference off-chain,
preserving full post-hoc auditability without on-chain
computational overhead.
LUNA encodes posteriors in fixed-point form under a global
scale factor, manages claims through an explicit on-chain
state machine supporting incremental out-of-order evidence
arrival, and provides a complete reconstructible inference
transcript from ledger state alone.
We implement LUNA on Ethereum-compatible infrastructure,
evaluate it on the Sepolia public testnet across BN
configurations of 4 to 32 evidence nodes, and compare
directly against full on-chain BN inference.
Results show that per-query gas is independent of BN size
and \placeholder{X}\% lower than on-chain inference at
32~nodes, while initialization cost scales linearly in the
number of CPT entries.
Fixed-point encoding under $\texttt{SCALE}=10^6$ introduces
a maximum absolute error of $5{\times}10^{-7}$ per posterior
with no decision inconsistencies across all tested
configurations.
Applied to physician-visit verification with four binary
evidence sources, LUNA demonstrates that interpretable,
auditable BN-based oracle decisions are practically
deployable in healthcare settings without prohibitive
on-chain costs, while identifying oracle centralization
and clinical CPT validation as the primary directions for
future work.
\end{abstract}

\begin{IEEEkeywords}
Bayesian inference, blockchain oracle, smart contracts,
healthcare informatics, gas optimization, evidence-based
reasoning, fixed-point arithmetic, claim verification
\end{IEEEkeywords}

%% ============================================================
\section{Introduction}
\label{sec:introduction}
%% ============================================================

\subsection{Motivation}

Smart contracts execute deterministically, yet their
correctness often depends on external data that contracts
cannot access directly.
As blockchain-based smart contracts increasingly mediate
high-stakes decisions in finance, supply chains, and
healthcare~\cite{hameed2024digital}, they rely on
\emph{oracles} to supply off-chain
information~\cite{sheldon2021auditing}.
This creates the \emph{oracle problem}: even with secure
on-chain consensus, contract correctness ultimately depends
on the trustworthiness and interpretability of off-chain
data feeds~\cite{pasdar2023connect}.

Most existing oracle designs address this by reducing
heterogeneous evidence to a scalar value: a price
feed~\cite{zhang2024blockchain,ahmadjee2025decision}, a
reputation
score~\cite{chen2022decentralized,nelaturu2020public,arshad2022reputable,
wang2025efficient}, or a majority
vote~\cite{adler2018astraea,sober2021voting}.
While sufficient for simple aggregation, scalar abstractions
obscure how individual evidence sources contribute to a
decision, making outcomes difficult to audit, contest, or
formally analyse~\cite{pasdar2023connect,ezzat2022blockchain}.
This limitation is especially consequential in healthcare,
where clinical and administrative decisions must
demonstrably follow from identifiable evidence, and where
regulatory frameworks require explainable, auditable
reasoning~\addcite.
%% TODO: Replace \addcite with a real citation, e.g.:
%% EU MDR 2017/745 (Article 61), IEEE 7001-2021, or a
%% JBHI paper on clinical decision support accountability.

Bayesian Networks~(BNs) offer a principled alternative.
By representing claims and evidence as nodes and encoding
probabilistic dependencies via Conditional Probability
Tables~(CPTs), BNs enable structured, interpretable
inference in which each evidence source has an explicit
and locally verifiable effect on the posterior.
Prior work has demonstrated that BN-based reasoning can be
embedded directly in smart contracts, allowing on-chain
posterior computation from evidence~\cite{bn-oracle}.
This approach has two concrete advantages over scalar
oracles: decisions are derived from an explicit
probabilistic model whose parameters are visible on-chain,
and the contribution of each individual evidence source is
traceable through its CPT entries.

However, on-chain BN inference faces a fundamental
scalability barrier.
CPT evaluation and posterior normalization require
on-chain operations whose gas cost grows with the number
of evidence nodes, making full on-chain BN inference
impractical for models with more than a small number of
sources~\cite{bn-oracle}.
The obvious alternative---moving inference entirely
off-chain---reduces gas consumption but severs the
verifiability link: without a principled interface, smart
contracts cannot determine whether a reported posterior is
consistent with the BN parameters anchored to the ledger.
This creates a trust gap that existing designs do not
close.

The core tension is therefore this: on-chain BN inference
preserves verifiability but does not scale; off-chain
inference scales but abandons verifiability.
Neither extreme is acceptable for healthcare settings,
where both computational feasibility and decision
auditability are non-negotiable requirements.
A further complication is oracle centralization: if a
single operator controls inference submission, correctness
is assumed rather than enforced at execution time.
LUNA addresses this partially through on-chain parameter
anchoring and post-hoc auditability, but full
decentralization of the inference layer remains an open
problem scoped as future work.

These considerations motivate the following research
questions:

\begin{itemize}
  \item \textbf{RQ1:} Can a BN-based oracle be designed
  such that per-query on-chain gas cost is independent of
  BN size, while BN parameters, evidence assignments, and
  posterior outputs remain anchored and auditable on-chain?

  \item \textbf{RQ2:} Can a realistic claim--evidence
  lifecycle---in which claims are opened once and updated
  incrementally as evidence arrives out of order---be
  supported without contract redeployment or modification
  of the BN structure?

  \item \textbf{RQ3:} Does off-chain inference with
  on-chain parameter anchoring preserve sufficient
  numerical fidelity under fixed-point encoding that
  posteriors remain decision-consistent with their
  floating-point counterparts, and what are the conditions
  under which this guarantee holds?
\end{itemize}

\subsection{Contributions}

This paper introduces \emph{LUNA} (Ledger-grounded
Uncertainty-aware Networked Architecture), a BN-based
oracle architecture that resolves the
scalability--verifiability tension by treating the
blockchain as an authenticated parameter and audit plane
rather than an inference engine.
Inference executes off-chain against a BN whose structure
and CPTs are anchored on-chain; only compact evidence
assignments and scaled posteriors are committed per query.
This keeps the on-chain interface fixed-width and
gas-predictable regardless of BN size, while preserving
the ability to reconstruct and independently verify every
inference decision from on-chain state.

The paper makes the following concrete contributions:

\begin{itemize}

  \item \textbf{Gas-predictable BN oracle architecture.}
  We design an oracle in which BN priors, CPTs, evidence
  assignments, and posterior outputs are anchored on-chain
  while inference executes off-chain.
  The on-chain interface is fixed-width by construction,
  yielding per-query gas cost that is independent of BN
  size under a fixed query model.
  We confirm this empirically across BN configurations
  spanning 4 to 32 evidence nodes on both a local
  development network and the Sepolia public testnet,
  and compare directly against full on-chain BN
  inference~\cite{bn-oracle} to quantify the gas
  reduction.

  \item \textbf{Formal evidence-based oracle model.}
  We define a restricted discrete BN with binary claim
  and evidence variables, derive a closed-form inference
  procedure, and specify a fixed-point encoding scheme.
  We characterise the rounding error introduced by
  fixed-point representation as a function of the scale
  factor and establish the conditions under which encoded
  posteriors remain decision-consistent with their
  floating-point counterparts.

  \item \textbf{Claim lifecycle and interaction protocol.}
  We specify an on-chain state machine supporting
  incremental, out-of-order evidence arrival.
  Claims are opened once and updated arbitrarily without
  redeployment.
  We show empirically that per-update gas cost is stable
  across the full claim lifecycle under multiple
  concurrent claims and evidence orderings.

  \item \textbf{Implementation and empirical evaluation.}
  We implement LUNA on Ethereum-compatible infrastructure
  and evaluate on the Sepolia public testnet and a local
  development network.
  Results show per-query gas is effectively constant as
  BN size grows, initialization cost scales linearly in
  CPT entries, and posteriors remain decision-consistent
  with off-chain computation across all tested
  configurations.
  CPT values are synthetic in this evaluation; clinical
  grounding of parameters is identified as the primary
  direction for future work.

\end{itemize}

Taken together, these contributions demonstrate that
interpretable, auditable BN-based oracle decisions are
practically deployable without gas costs that scale with
model complexity---a combination not achieved by prior
designs.
The architecture is explicit about its current
limitation: auditability is post-hoc, and the inference
layer remains centralized pending future work on
decentralized oracle committees.

The remainder of the paper is organized as follows.
Section~\ref{sec:relatedwork} reviews related work.
Section~\ref{sec:proposed-methodology} presents the LUNA
system model and BN formulation.
Section~\ref{sec:implementation} describes the
implementation.
Section~\ref{sec:experimental-analysis} reports
experimental results.
Section~\ref{sec:results-discussion} discusses trade-offs
and limitations.
Section~\ref{sec:conclusion} concludes.


%% ============================================================
\section{Related Work}
\label{sec:relatedwork}
%% ============================================================

Blockchain oracles provide external data to smart
contracts but differ significantly in how they aggregate
evidence, represent uncertainty, and ensure trust.
We survey four categories relevant to LUNA and summarize
key distinctions in Table~\ref{tab:oracle-comparison}.

\subsection{Data-Feed and Oracle Networks}

Widely deployed oracle systems such as
Chainlink~\cite{breidenbach2021chainlink} and
Witnet~\cite{de2017witnet} provide smart contracts
with scalar values---asset prices or event
outcomes---sourced from multiple off-chain providers.
These systems rely on aggregation mechanisms,
cryptographic signatures, and decentralized committees
to improve robustness against faulty or malicious data
providers~\cite{pasdar2023connect,ezzat2022blockchain,mammadzada2020blockchain}.
While effective for numerical feeds, such designs treat
inputs as atomic values and do not explicitly model how
different evidence sources contribute to a decision.
The reasoning process behind a reported value remains
implicit, limiting interpretability and making it
difficult to audit how individual signals affect
outcomes.

\subsection{Reputation, Voting, and Prediction-Market
Oracles}

Alternative oracle designs use economic incentives and
collective decision-making.
Reputation-based systems assign trust scores to data
providers~\cite{chen2022decentralized,nelaturu2020public,
arshad2022reputable}, while voting-based and
prediction-market oracles such as
Augur~\cite{peterson2015augur} aggregate reports through
token-weighted voting and reward mechanisms.
Kleros~\cite{lesaege2019kleros} extends this with
cryptoeconomic dispute resolution.
These approaches provide economic security and
decentralization, but compress heterogeneous evidence
into scalar scores or discrete outcomes.
The relationship between individual evidence sources and
final decisions is not explicitly represented, limiting
transparency and making post-hoc analysis of decisions
difficult~\cite{pasdar2021blockchain}.

\subsection{Verifiable Computation and zk-Oracles}

A separate line of work focuses on verifiable
computation, in which off-chain results are accompanied
by cryptographic proofs verified on-chain.
Systems based on trusted execution environments such as
Town~Crier~\cite{zhang2016town} or zero-knowledge
proofs such as DECO~\cite{zhang2020deco} and
zk-SNARK-based oracles~\cite{garg2022succinct} ensure
that a reported output is consistent with a specified
computation.
These approaches provide strong correctness guarantees
but typically treat the underlying computation as a
black box.
While they can attest that a result was computed
correctly, they do not prescribe how to structure the
model itself or how to expose intermediate evidence and
parameters for interpretability.
Proof generation and verification also introduce
non-trivial overhead that scales with circuit complexity.

\subsection{Bayesian Network--Based Oracles}

Recent work explores the use of BNs to represent
structured evidence in oracle systems~\cite{bn-oracle}.
In these designs, claims and evidence are modeled
explicitly, and smart contracts compute posterior
probabilities directly from on-chain CPTs.
This enables interpretable, evidence-based decisions
with fully transparent parameters.
However, executing BN inference on-chain incurs gas
costs that grow with the size of the network and the
number of evidence nodes, making such approaches
difficult to scale beyond small models.
Prior BN-based systems also remain largely
domain-agnostic decision aids; they do not address how
to commit BN parameters, evidence configurations, and
resulting decisions to a permissionless ledger in a
gas-efficient yet auditable manner.

\subsection{Positioning LUNA}

LUNA occupies a distinct position in this landscape.
Unlike data-feed oracles, LUNA makes the probabilistic
model explicit: each evidence source is a node in a BN
whose CPTs are publicly anchored on-chain.
Unlike reputation and voting systems, LUNA produces
interpretable posteriors with a clear semantic link to
individual evidence sources.
Unlike verifiable computation approaches, LUNA does not
require proof generation; instead, it achieves
auditability through parameter transparency and
transcript completeness, trading runtime correctness
enforcement for architectural simplicity.
Unlike prior BN-on-chain designs, LUNA decouples
inference from the on-chain interface, achieving
gas-predictable operation regardless of BN size.
Table~\ref{tab:oracle-comparison} summarizes these
distinctions.

\begin{table}[t]
\centering
\caption{Comparison of oracle designs along inference,
accountability, and gas cost dimensions.}
\label{tab:oracle-comparison}
\resizebox{\columnwidth}{!}{%
\begin{tabular}{p{2.0cm}p{1.3cm}p{2.9cm}p{2.0cm}}
\toprule
\textbf{Category} &
\textbf{Inference} &
\textbf{Accountability} &
\textbf{Gas cost} \\
\midrule
Price / data-feed
oracles~\cite{breidenbach2021chainlink,
de2017witnet}
& Hybrid
& Values, signatures, committee votes
& Tied to feed size; model implicit \\
\addlinespace
Reputation / prediction-market
oracles~\cite{peterson2015augur,lesaege2019kleros}
& Off-chain / hybrid
& Reputation scores, staking, slashing, juror voting
& Depends on number of reporters \\
\addlinespace
Verifiable-computation
oracles~\cite{zhang2016town,zhang2020deco,
garg2022succinct}
& Off-chain
& Proofs on-chain; computation verifiable
& Proof generation cost \\
\addlinespace
BN-on-chain
oracles~\cite{bn-oracle}
& On-chain
& Parameters and posteriors fully on-chain
& Grows with BN size \\
\addlinespace
\textbf{LUNA (this work)}
& Off-chain
& Priors, CPTs, claim states, and per-claim
  posteriors recorded on-chain
& Per-query gas $\approx$ constant \\
\bottomrule
\end{tabular}}
\end{table}


%% ============================================================
\section{System Model and Architecture}
\label{sec:proposed-methodology}
%% ============================================================

LUNA treats the blockchain as an authenticated parameter
and audit plane rather than an inference engine.
A discrete BN encodes the probabilistic relationship
between claims and evidence; inference executes off-chain
against BN parameters anchored on-chain; and every
evidence assignment, posterior output, and state
transition is committed to the ledger as an immutable
transcript.
This separation achieves gas-predictable on-chain
interaction while preserving the ability to reconstruct
and independently verify every inference decision from
ledger state alone.
Figure~\ref{fig:component} shows the component structure;
Figure~\ref{fig:sequence} shows the interaction protocol.

\subsection{Actors and Roles}

LUNA involves four roles.

\begin{itemize}

  \item \textbf{BN Designer.}
  Specifies BN structure, node semantics, and initial
  CPT parameters. Deploys on-chain contracts and
  initialises \texttt{CPTStore}. Trusted for model
  correctness; CPT values are publicly visible on-chain
  and independently auditable by any party.

  \item \textbf{Oracle Operator.}
  Runs the off-chain BN inference service. Reconstructs
  BN parameters from \texttt{CPTStore}, evaluates
  posteriors for incoming evidence assignments, and
  submits results through \texttt{OracleController}.
  In the current design this is a single trusted party;
  correctness of submitted posteriors is auditable
  post-hoc but is not enforced at submission time.
  Decentralisation of this role is identified as future
  work (Section~\ref{sec:limitations}).

  \item \textbf{Evidence Provider.}
  Submits binary evidence observations linked to a
  specific claim. Each submission is
  signature-verified on-chain. Evidence providers are
  assumed honest in the current model; adversarial
  evidence injection is not defended against at the
  contract level.

  \item \textbf{Querier.}
  Reads committed posteriors and audit logs from the
  ledger. Requires no special privileges; all committed
  state is publicly readable.

\end{itemize}

The system enforces three architectural invariants by
construction. First, \textbf{structural fixity}: the BN
topology, node semantics, and admissible query type are
fixed at deployment time and cannot be modified without
redeployment, ensuring all claims under a given
deployment are evaluated against an identical model.
Second, \textbf{inference separation}: smart contracts
never execute probabilistic inference; the on-chain
interface accepts only fixed-width tuples of evidence
bits and scaled posterior values, making per-query gas
independent of BN structure.
Third, \textbf{transcript completeness}: every parameter
write, evidence submission, and posterior commitment is
recorded in contract storage or emitted as an indexed
event, making the full reasoning transcript for any
claim reconstructible from ledger state alone.

\subsection{Threat Model}

LUNA inherits the tamper-evidence and append-only
guarantees of the underlying blockchain for all
on-chain state. Within this boundary, four threats
are in scope.

\textbf{Malicious posterior submission.}
The oracle operator could submit posteriors inconsistent
with the anchored CPTs. LUNA does not prevent this at
execution time. Mitigation is post-hoc: any party can
reconstruct the correct posterior from \texttt{CPTStore}
and the committed evidence assignment and detect a
discrepancy. Future work will address this through
decentralised oracle committees or verifiable inference
using zero-knowledge proofs.

\textbf{Stale parameter submission.}
CPT parameters may be updated between an oracle's
reconstruction and its posterior submission, causing
the submitted posterior to reflect outdated parameters.
LUNA mitigates this by binding each posterior submission
to the block height at which parameters were
reconstructed, allowing auditors to detect
parameter-submission races.

\textbf{Adversarial evidence injection.}
A compromised evidence provider could submit false
binary observations. On-chain signature verification
ensures only authorised providers can submit evidence,
but correctness of the observation itself is not
verified. This is a data quality assumption common to
oracle systems and is noted as a limitation.

\textbf{Out-of-scope threats.}
Smart contract vulnerabilities (re-entrancy, integer
overflow) are mitigated by Solidity~0.8.x checked
arithmetic and are not analysed further.
Blockchain-level attacks (51\% attack, long-range
attacks) are assumed negligible under the target
deployment conditions.

\subsection{Bayesian Network Formulation}

LUNA uses a discrete BN over two binary claim
variables, $\texttt{PPH}$~(in-person physician visit)
and $\texttt{PPR}$~(remote physician visit), and four
binary evidence variables
\[
  \mathcal{E}=\{
    \texttt{GPS},\,\texttt{PC},\,
    \texttt{PMD},\,\texttt{PR}
  \},
\]
where \texttt{GPS} denotes GPS location consistency,
\texttt{PC} patient confirmation, \texttt{PMD}
physician medical device log, and \texttt{PR}
physician prescription record.
Each evidence variable is conditioned on both claim
variables:
\[
  (\texttt{PPH},\texttt{PPR})
  \;\rightarrow\;
  \{\texttt{GPS},\texttt{PC},\texttt{PMD},\texttt{PR}\}.
\]

\noindent\textbf{Conditional independence assumption.}
Given the claim configuration $(\texttt{PPH},
\texttt{PPR})$, evidence variables are modeled as
conditionally independent, yielding the factorisation
\begin{equation}
  P(\texttt{PPH},\texttt{PPR},E)
  =
  P(\texttt{PPH})\,P(\texttt{PPR})
  \prod_{E_i\in\mathcal{E}}
  P(E_i\mid \texttt{PPH},\texttt{PPR}).
  \label{eq:factorization}
\end{equation}
This is a strong assumption: in practice, GPS
consistency and patient confirmation may be correlated
independently of the claim type. We adopt it to keep
inference closed-form and constant-cost; relaxing it
to a richer dependency structure is left to future
work. Sensitivity of posterior outputs to this
assumption is evaluated in
Section~\ref{sec:exp-sensitivity}.

\noindent\textbf{Claim variable independence.}
\texttt{PPH} and \texttt{PPR} are modeled as
marginally independent root nodes. This does not imply
that in-person and remote visits can co-occur; mutual
exclusivity is enforced at the application layer by
the claim submission protocol. An alternative encoding
as a single binary root variable
$C\in\{\texttt{PPH},\texttt{PPR}\}$ would simplify the
structure and is considered in
Section~\ref{sec:results-discussion}.

\noindent\textbf{Posterior computation.}
For a given evidence assignment
$e=(e_1,\dots,e_4)\in\{0,1\}^4$, the oracle evaluates
\begin{equation}
  P(\texttt{PPH},\texttt{PPR}\mid e)
  \;\propto\;
  P(\texttt{PPH})\,P(\texttt{PPR})
  \prod_{E_i\in\mathcal{E}}
  P(E_i=e_i\mid \texttt{PPH},\texttt{PPR}),
  \label{eq:posterior}
\end{equation}
enumerating all four root configurations
$(i,j)\in\{0,1\}^2$ and normalising.
The committed marginals $P(\texttt{PPH}=1\mid e)$ and
$P(\texttt{PPR}=1\mid e)$ are obtained by summing over
the complementary root value. Inference complexity is
$O(2^{N_h}\cdot N_e)$ for $N_h$ root variables and
$N_e$ evidence nodes, reducing to a constant-time
computation for fixed $N_h=2$.

\subsection{Fixed-Point Encoding and Decision
Consistency}

Posterior probabilities $p\in[0,1]$ are encoded as
integers using a global scale factor
$\texttt{SCALE}\in\mathbb{Z}^+$:
\[
  \widehat{p}
  = \texttt{round}(p\cdot\texttt{SCALE}),
  \qquad
  \widehat{p}\in\{0,\dots,\texttt{SCALE}\}.
\]

\begin{proposition}[Rounding error bound]
The absolute rounding error satisfies
\[
  \left|p - \frac{\widehat{p}}{\texttt{SCALE}}\right|
  \leq \frac{1}{2\cdot\texttt{SCALE}}.
\]
For $\texttt{SCALE}=10^6$, the per-posterior error is
bounded by $5\times10^{-7}$.
\end{proposition}

\begin{definition}[Decision consistency]
A posterior $p$ is \emph{decision-consistent} with its
encoded counterpart $\widehat{p}$ at threshold $\tau$
if and only if
\[
  \mathbb{1}[p \geq \tau]
  \;=\;
  \mathbb{1}\!\left[
    \tfrac{\widehat{p}}{\texttt{SCALE}} \geq \tau
  \right].
\]
Inconsistency is possible only when
$p\in[\tau-\epsilon,\,\tau+\epsilon]$ where
$\epsilon = \frac{1}{2\cdot\texttt{SCALE}}$.
\end{definition}

For $\texttt{SCALE}=10^6$ and $\tau=0.5$, the
inconsistency interval has width $10^{-6}$.
Whether any evaluated posterior falls in this interval
is reported empirically in
Section~\ref{sec:exp-fidelity}.

\subsection{Claim Lifecycle and Oracle Workflow}

Each claim represents a single physician visit,
identified by an internal on-chain identifier
\texttt{claimId} and an external key
$k\in\{0,1\}^{256}$ derived from visit metadata via
\texttt{keccak256}.
Evidence may arrive incrementally and out of order.
The current state of a claim at time $t$ is a partial
assignment
\[
  e^{(t)}:
  \{\texttt{GPS},\texttt{PC},\texttt{PMD},\texttt{PR}\}
  \rightarrow\{0,1,\bot\},
\]
where $\bot$ denotes a variable not yet observed.
Inference is computed over the observed subset only;
unobserved variables are marginalised out.

Each claim transitions through three states:
\[
  \texttt{None}
  \xrightarrow{\texttt{openClaim}}
  \texttt{Open}
  \xrightarrow{\texttt{finalizeClaim}}
  \texttt{Finalized}.
\]
Posterior submissions via \texttt{submitInference} are
accepted only in the \texttt{Open} state.
The \texttt{Finalized} state is absorbing; no further
updates are possible, preventing retroactive
modification of completed claims.

The end-to-end workflow proceeds as follows.
(i)~An application calls $\texttt{openClaim}(k)$
on-chain, creating a new claim in the \texttt{Open}
state.
(ii)~Evidence providers submit binary observations
linked to \texttt{claimId}; each is signature-verified
and appended to \texttt{EvidenceRegistry}.
(iii)~The off-chain oracle reconstructs $e^{(t)}$
from \texttt{EvidenceRegistry}, retrieves BN
parameters from \texttt{CPTStore} at block~$b$,
evaluates Equation~\eqref{eq:posterior}, encodes both
marginals, and calls \texttt{submitInference}.
(iv)~\texttt{OracleController} validates the
submission, updates \texttt{ClaimRegistry}, appends to
\texttt{EvidenceRegistry}, and emits an indexed event
binding the posterior to block height~$b$.
(v)~Once complete, \texttt{finalizeClaim} transitions
the claim to \texttt{Finalized}.

Because every parameter write, evidence submission, and
posterior commit is recorded on-chain, the full
inference transcript for any claim is reconstructible
and auditable from ledger state alone without reliance
on off-chain records.

\begin{figure}[t]
  \centering
  \IfFileExists{Figures/Component_image.png}%
  {\includegraphics[width=\columnwidth]{Figures/Component_image}}%
  {\fbox{\parbox{\columnwidth}{\centering\vspace{1.5cm}
   \texttt{[insert: Figures/Component\_image.png]}
   \vspace{1.5cm}}}}
  \caption{LUNA component architecture. The on-chain
  layer hosts four smart contracts; the off-chain oracle
  performs BN inference and submits posteriors through
  \texttt{OracleController}.}
  \label{fig:component}
\end{figure}

\begin{figure}[t]
  \centering
  \IfFileExists{Figures/sequence_diagram.png}%
  {\includegraphics[width=\columnwidth]{Figures/sequence_diagram}}%
  {\fbox{\parbox{\columnwidth}{\centering\vspace{1.5cm}
   \texttt{[insert: Figures/sequence\_diagram.png]}
   \vspace{1.5cm}}}}
  \caption{LUNA interaction protocol covering BN
  creation and approval, claim assertion, evidence
  assertion, and claim verification phases.}
  \label{fig:sequence}
\end{figure}


%% ============================================================
\section{Implementation}
\label{sec:implementation}
%% ============================================================

The complete implementation---smart contracts,
off-chain oracle, and experiment scripts---is publicly
available at
\url{https://github.com/S1120957/S_Evidence_BN_Oracle}
(archived: DOI~\placeholder{INSERT ZENODO DOI},
tag~\texttt{v1.0}).

\subsection{On-Chain Contracts}

The on-chain layer is implemented in Solidity~0.8.x
on an Ethereum-compatible EVM, organized into four
contracts with strictly separated responsibilities.
Table~\ref{tab:access-control} summarizes write
permissions.

\textbf{CPTStore} is the sole repository for BN
parameters. It stores prior probabilities
$P(\texttt{PPH}=1)$ and $P(\texttt{PPR}=1)$, and for
each evidence node $E_i\in\mathcal{E}$, the four
conditional entries
$P(E_i=1\mid\texttt{PPH}=i,\texttt{PPR}=j)$,
$i,j\in\{0,1\}$. All values are stored as integers
$\widehat{p}\in\{0,\dots,\texttt{SCALE}\}$ under
$\texttt{SCALE}=10^6$. Every parameter write is
range-checked and emitted as an indexed event,
enabling full reconstruction of the BN parameter
history from the ledger. Write access is restricted
to the deploying account; reads are unrestricted.

\textbf{ClaimRegistry} manages claim state. Each
claim is identified by an internal \texttt{claimId}
and linked to an external key via \texttt{keccak256}.
The registry stores the latest committed posterior
pair and enforces the state machine
$\texttt{None}\to\texttt{Open}\to\texttt{Finalized}$.
The \texttt{Finalized} state is absorbing. Only
\texttt{OracleController} may write to this registry.

\textbf{EvidenceRegistry} maintains an append-only
log of evidence--posterior pairs. Each entry records
the \texttt{claimId}, four evidence bits, scaled
posteriors, submitting address, and block number.
Given this log and the parameter history in
\texttt{CPTStore}, any party can independently
recompute and verify every posterior ever committed
for any claim. Write-restricted to
\texttt{OracleController}.

\textbf{OracleController} is the sole external entry
point for inference submission. It accepts a
fixed-width tuple
$(\texttt{claimId},\,e\in\{0,1\}^4,\,
\widehat{p}_{\texttt{PPH}},\,
\widehat{p}_{\texttt{PPR}})$, enforces three
invariants: (i)~claim exists and is \texttt{Open},
(ii)~posterior values satisfy
$0\leq\widehat{p}\leq\texttt{SCALE}$,
(iii)~caller is the authorised oracle address.
On success it updates \texttt{ClaimRegistry}, appends
to \texttt{EvidenceRegistry}, and emits an indexed
event binding the submission to CPTStore block
height~$b$. No contract reads CPT entries or performs
arithmetic on posterior values; all operations are
$O(1)$ in BN size by construction.

\begin{table}[t]
\centering
\caption{Write permission matrix for LUNA contracts.
Integer overflow is precluded by Solidity~0.8.x
checked arithmetic throughout.}
\label{tab:access-control}
\begin{tabularx}{\columnwidth}{lXl}
\toprule
\textbf{Contract} &
\textbf{Written data} &
\textbf{Caller} \\
\midrule
\texttt{CPTStore}
  & BN priors and CPT entries
  & Deployer \\
\texttt{ClaimRegistry}
  & Claim state and posteriors
  & \texttt{OracleController} \\
\texttt{EvidenceRegistry}
  & Evidence--posterior log
  & \texttt{OracleController} \\
\texttt{OracleController}
  & Inference submissions
  & Authorised oracle \\
\bottomrule
\end{tabularx}
\end{table}

\subsection{Off-Chain Oracle}

The off-chain oracle is implemented in Python~3.x.
Its core component, \texttt{BNOracle}, stores BN
structure and floating-point CPT parameters and
evaluates the posterior pair
$(P(\texttt{PPH}=1\mid e),\,P(\texttt{PPR}=1\mid e))$
via Equation~\eqref{eq:posterior}.

Before each inference run, a synchronisation module
reads all CPT entries from \texttt{CPTStore} via
\texttt{web3.py} and reconstructs the BN from the
current on-chain parameter snapshot at block~$b$.
This ensures submitted posteriors correspond to the
parameter state visible on-chain at reconstruction
time. The block height~$b$ is recorded and included
in the submission event for audit purposes. If
reconstruction fails (RPC timeout or unexpected
contract state), the oracle aborts and logs the
failure without writing to the chain.

The submission pipeline executes five steps:
(i)~reconstruct BN from \texttt{CPTStore} at block~$b$;
(ii)~read current evidence assignment $e^{(t)}$ from
\texttt{EvidenceRegistry};
(iii)~evaluate Equation~\eqref{eq:posterior} over
observed variables, marginalising out $\bot$ entries;
(iv)~encode both marginals as
$\widehat{p}=\texttt{round}(p\cdot\texttt{SCALE})$;
(v)~submit via \texttt{submitInference}, recording
transaction hash and gas used.

Transaction reverts are caught and logged with the
revert reason. The oracle does not retry automatically;
failed submissions require manual inspection. This is
a known limitation of the single-operator design
discussed in Section~\ref{sec:limitations}.

\subsection{Deployment and Experimental Setup}

\textbf{Deployment environments.}
Experiments use two complementary environments.
\emph{Ganache} provides a local deterministic network
with fixed gas price~(20~Gwei), isolating pure EVM
execution cost from network variability.
\emph{Sepolia} is the Ethereum public testnet
(EIP-1559 base fees, variable block packing), serving
as the primary validity check for the
gas-predictability claim. All Sepolia transactions are
publicly verifiable at the addresses in
Table~\ref{tab:contracts}.

\begin{table}[t]
\centering
\caption{Deployed contract addresses on the Sepolia
public testnet, verifiable via Sepolia Etherscan.}
\label{tab:contracts}
\begin{tabularx}{\columnwidth}{lX}
\toprule
\textbf{Contract} & \textbf{Sepolia address} \\
\midrule
\texttt{CPTStore}
  & \placeholder{0x...} \\
\texttt{ClaimRegistry}
  & \placeholder{0x...} \\
\texttt{EvidenceRegistry}
  & \placeholder{0x...} \\
\texttt{OracleController}
  & \placeholder{0x...} \\
\bottomrule
\end{tabularx}
\end{table}

\textbf{Toolchain.}
Contracts were compiled and deployed using
Hardhat~2.x. Python components interacted with both
networks via \texttt{web3.py}~6.x. All dependencies
are version-pinned in \texttt{requirements.txt} and
\texttt{package.json}.

\textbf{BN parameter configurations.}
Two CPT configurations are evaluated. The
\emph{neutral} configuration sets all priors and CPT
entries to $p=0.5$ ($\widehat{p}=500{,}000$),
isolating architectural properties from
parameter-specific effects. The \emph{asymmetric}
configuration uses values specified in
Table~\ref{tab:cpt-values}, motivated by
plausible clinical priors~\addcite.
%% TODO: Add citation justifying asymmetric CPT values.
Gas behaviour and fixed-point fidelity are evaluated
under both configurations. Neutral CPTs produce
uninformative posteriors by construction; posterior
convergence results are therefore reported only for
the asymmetric configuration.

\begin{table}[t]
\centering
\caption{Asymmetric CPT configuration.
Values are $P(E_i=1\mid\texttt{PPH}=i,\texttt{PPR}=j)$.
\todo{Justify values with a clinical or domain
reference before submission.}}
\label{tab:cpt-values}
\resizebox{\columnwidth}{!}{%
\begin{tabular}{lcccc}
\toprule
& \multicolumn{4}{c}{\textbf{Parent configuration}
  $(\texttt{PPH},\texttt{PPR})$} \\
\cmidrule{2-5}
\textbf{Evidence} &
$(1,0)$ & $(0,1)$ & $(1,1)$ & $(0,0)$ \\
\midrule
\texttt{GPS} &
\placeholder{0.90} & \placeholder{0.15} &
\placeholder{0.80} & \placeholder{0.10} \\
\texttt{PC}  &
\placeholder{0.85} & \placeholder{0.20} &
\placeholder{0.75} & \placeholder{0.15} \\
\texttt{PMD} &
\placeholder{0.88} & \placeholder{0.10} &
\placeholder{0.78} & \placeholder{0.08} \\
\texttt{PR}  &
\placeholder{0.80} & \placeholder{0.25} &
\placeholder{0.70} & \placeholder{0.20} \\
\bottomrule
\end{tabular}}
\end{table}

\textbf{Experiment profiles.}
Three BN sizes are evaluated: 4~nodes (baseline),
8~nodes, and 32~nodes. Evidence nodes beyond the
baseline four are synthetic, included only in the
off-chain BN and \texttt{CPTStore}; the
\texttt{OracleController} interface remains
fixed-width across all sizes. For each size and CPT
configuration, all $2^4=16$ complete evidence
assignments are evaluated, plus partial assignments
simulating incremental out-of-order arrival across
three concurrent claims with different orderings.
On Sepolia, each transaction type is submitted
$n=20$ times; mean and standard deviation are
reported.


%% ============================================================
\section{Experimental Evaluation}
\label{sec:experimental-analysis}
%% ============================================================

We evaluate LUNA along four dimensions:
(i)~on-chain gas cost as a function of BN size,
compared directly against full on-chain BN inference;
(ii)~fixed-point fidelity and decision consistency;
(iii)~claim lifecycle stability under incremental
out-of-order evidence arrival; and
(iv)~end-to-end pipeline feasibility on a public
testnet.

\subsection{Gas Scaling with BN Size}
\label{sec:exp-bn-size}

Table~\ref{tab:gas-scaling} reports gas consumption
for three transaction types across BN sizes on Ganache
and Sepolia.

\begin{table}[t]
\centering
\caption{Gas consumption by transaction type and BN
size. Sepolia values: mean~$\pm$~std over $n=20$
independent submissions.
\todo{Fill all cells with measured values.}}
\label{tab:gas-scaling}
\resizebox{\columnwidth}{!}{%
\begin{tabular}{llrrr}
\toprule
\textbf{Network} &
\textbf{Transaction} &
\textbf{4 nodes} &
\textbf{8 nodes} &
\textbf{32 nodes} \\
\midrule
Ganache
  & \texttt{CPTStore} init
    & \placeholder{--}
    & \placeholder{--}
    & \placeholder{--} \\
  & \texttt{openClaim}
    & \placeholder{--}
    & \placeholder{--}
    & \placeholder{--} \\
  & \texttt{submitInference}
    & \placeholder{--}
    & \placeholder{--}
    & \placeholder{--} \\
\midrule
Sepolia
  & \texttt{CPTStore} init
    & \placeholder{--$\pm$--}
    & \placeholder{--$\pm$--}
    & \placeholder{--$\pm$--} \\
  & \texttt{openClaim}
    & \placeholder{--$\pm$--}
    & \placeholder{--$\pm$--}
    & \placeholder{--$\pm$--} \\
  & \texttt{submitInference}
    & \placeholder{--$\pm$--}
    & \placeholder{--$\pm$--}
    & \placeholder{--$\pm$--} \\
\bottomrule
\end{tabular}}
\end{table}

\textbf{Initialization cost.}
\texttt{CPTStore} initialization cost grows linearly
with the number of evidence nodes, as each CPT entry
requires one constant-cost \texttt{SSTORE} operation.
Measured costs are \placeholder{X}~gas at 4~nodes,
\placeholder{Y}~gas at 8~nodes, and \placeholder{Z}~gas
at 32~nodes, consistent with the $O(N_e)$ bound.

\textbf{Per-query cost.}
\texttt{submitInference} gas is
\placeholder{A$\pm$B}~gas on Sepolia across all BN
sizes, confirming that per-query cost is independent
of BN size and directly answering~RQ1. The on-chain
function performs a fixed number of \texttt{SSTORE}
and \texttt{LOG} operations regardless of BN
structure.

Figure~\ref{fig:gas-scaling} illustrates the flat
per-query gas profile against the linearly growing
initialization cost.

\begin{figure}[t]
  \centering
  %% TODO: Generate and insert gas scaling plot.
  %% x-axis: BN size (4, 8, 32 nodes)
  %% y-axis: gas units
  %% Two series: submitInference (flat line),
  %%             CPTStore init (rising line)
  %% Two panels or overlaid markers: Ganache and Sepolia
  \placeholder{[INSERT GAS SCALING FIGURE]}
  \caption{Gas consumption vs.\ BN size on Ganache and
  Sepolia. \texttt{submitInference} gas (solid) is
  effectively constant; \texttt{CPTStore}
  initialization (dashed) grows linearly in the number
  of CPT entries.}
  \label{fig:gas-scaling}
\end{figure}

\subsection{Comparison Against On-Chain BN Inference}
\label{sec:exp-comparison}

Table~\ref{tab:gas-comparison} compares LUNA's
per-query gas against the reference on-chain BN
inference contract~\cite{bn-oracle} across identical
BN configurations on Sepolia.

\begin{table}[t]
\centering
\caption{Per-query gas: LUNA \texttt{submitInference}
vs.\ on-chain BN inference~\cite{bn-oracle}.
\todo{Fill all cells with measured values.}}
\label{tab:gas-comparison}
\begin{tabularx}{\columnwidth}{lXXr}
\toprule
\textbf{BN size} &
\textbf{LUNA} &
\textbf{On-chain BN} &
\textbf{Reduction} \\
\midrule
4 nodes  &
  \placeholder{--} & \placeholder{--} &
  \placeholder{--\%} \\
8 nodes  &
  \placeholder{--} & \placeholder{--} &
  \placeholder{--\%} \\
32 nodes &
  \placeholder{--} & \placeholder{--} &
  \placeholder{--\%} \\
\bottomrule
\end{tabularx}
\end{table}

At 4~evidence nodes, LUNA requires \placeholder{A}~gas
versus \placeholder{C}~gas for on-chain inference, a
reduction of \placeholder{D}\%.
At 32~nodes, the gap widens to \placeholder{E}~versus
\placeholder{F}~gas (\placeholder{G}\%~reduction),
confirming that the gas advantage of LUNA's design
grows with BN size.

\subsection{Fixed-Point Fidelity and Decision
Consistency}
\label{sec:exp-fidelity}

We evaluate fixed-point fidelity across all $2^4=16$
complete evidence assignments under both neutral and
asymmetric CPT configurations. For each assignment
we compute
$\delta=|p-\widehat{p}/\texttt{SCALE}|$ and check
decision consistency at $\tau=0.5$.
We additionally test
$\texttt{SCALE}\in\{10^2,10^3,10^4,10^6\}$ to
identify the minimum scale sufficient for decision
consistency.

\begin{table}[t]
\centering
\caption{Fixed-point fidelity under
$\texttt{SCALE}=10^6$.
Near-boundary counts posteriors within
$\epsilon=5{\times}10^{-7}$ of $\tau=0.5$.
\todo{Fill all cells with measured values.}}
\label{tab:fidelity}
\begin{tabularx}{\columnwidth}{lXXXX}
\toprule
\textbf{CPT config} &
$\delta_{\max}$ &
$\bar{\delta}$ &
\textbf{Near-bdry} &
\textbf{Inconsist.} \\
\midrule
Neutral    &
  \placeholder{--} & \placeholder{--} &
  \placeholder{--} & \placeholder{--} \\
Asymmetric &
  \placeholder{--} & \placeholder{--} &
  \placeholder{--} & \placeholder{--} \\
\bottomrule
\end{tabularx}
\end{table}

Maximum observed error is \placeholder{X}, within the
theoretical bound of $5\times10^{-7}$.
\placeholder{No / N} decision inconsistencies are
observed under either configuration for
$\texttt{SCALE}=10^6$, answering~RQ3.
The minimum sufficient scale for decision consistency
under the asymmetric configuration is
$\texttt{SCALE}\geq\placeholder{THRESHOLD}$, providing
a practical lower bound for storage-constrained
deployments.

\subsection{Claim Lifecycle Evaluation}
\label{sec:exp-lifecycle}

We evaluate lifecycle stability using three concurrent
claims with evidence arriving in different orderings
under the asymmetric CPT configuration.
Claim~1 receives evidence in canonical order
$\texttt{GPS}\to\texttt{PC}\to\texttt{PMD}\to
\texttt{PR}$; Claim~2 in reverse order; Claim~3 with
a gap (\texttt{GPS} and \texttt{PR} first, then
\texttt{PC}, then \texttt{PMD}).

\textbf{Gas stability.}
\texttt{submitInference} gas is
\placeholder{A$\pm$B}~gas across all stages and all
three claims, confirming that neither evidence
ordering nor partial observation affects on-chain cost
and directly answering~RQ2.

\textbf{Posterior convergence.}
Figure~\ref{fig:lifecycle} shows posterior
trajectories for all three claims. Under asymmetric
CPTs, posteriors for Claims~1 and~2 converge
toward \placeholder{PPH/PPR} as evidence accumulates,
demonstrating consistent convergence regardless of
arrival order. Claim~3 shows
\placeholder{describe actual behaviour} under partial
observation, confirming stability under missing
evidence.

\begin{figure}[t]
  \centering
  %% TODO: Generate and insert lifecycle figure.
  %% 3 subplots (Claims 1, 2, 3)
  %% x-axis: evidence stage 0 to 4
  %% y-axis: posterior probability [0,1]
  %% Two lines per plot: P(PPH=1|e^t) solid,
  %%                     P(PPR=1|e^t) dashed
  \placeholder{[INSERT LIFECYCLE POSTERIOR FIGURE]}
  \caption{Posterior evolution across evidence stages
  for three concurrent claims under asymmetric CPTs.
  $P(\texttt{PPH}=1\mid e^{(t)})$~(solid) and
  $P(\texttt{PPR}=1\mid e^{(t)})$~(dashed).}
  \label{fig:lifecycle}
\end{figure}

\textbf{State transitions.}
All three claims transition deterministically through
the state machine. Post-finalisation submission
attempts produce the expected revert in all three
cases, confirming that the \texttt{Finalized} state
is absorbing.

\subsection{Sensitivity Analysis}
\label{sec:exp-sensitivity}

\todo{Run sensitivity experiments before submission.
Suggested approach: vary each CPT entry across
\{0.1, 0.5, 0.9\} while holding others fixed; report
how posterior changes for a fixed evidence assignment;
confirm gas cost is invariant to CPT values. This
defends the conditional independence assumption in
Section~III and shows the architecture is robust to
parameter perturbation.}


%% ============================================================
\section{Discussion}
\label{sec:results-discussion}
%% ============================================================

\subsection{Answering the Research Questions}

\textbf{RQ1 --- Gas-independent per-query cost.}
The results confirm that per-query gas is independent
of BN size by construction, with
\placeholder{submitInference = A$\pm$B~gas} on Sepolia
across all tested sizes and a \placeholder{G}\%
reduction relative to on-chain BN inference at
32~nodes. Initialization is the only BN-size-dependent
cost and is paid once at deployment.

\textbf{RQ2 --- Lifecycle support without
redeployment.}
Three concurrent claims with different evidence
orderings all complete the full lifecycle without
redeployment, with stable per-update gas cost
independent of evidence history length or ordering.

\textbf{RQ3 --- Fixed-point decision consistency.}
For $\texttt{SCALE}=10^6$, no decision
inconsistencies are observed across all tested
configurations. The minimum sufficient scale is
$\texttt{SCALE}\geq\placeholder{THRESHOLD}$,
providing a practical lower bound for deployments
with tighter storage constraints.

\subsection{Design Trade-offs}

The core design choice---off-chain inference with
on-chain parameter anchoring---introduces a clear
trade-off between gas efficiency and runtime
correctness enforcement.

\textbf{Auditability vs.\ real-time enforcement.}
LUNA provides post-hoc auditability: any party can
reconstruct and verify any posterior from on-chain
state. However, a malicious oracle operator can
submit an incorrect posterior detectable only after
the fact. In healthcare settings where decisions are
acted upon immediately, post-hoc detection may be
insufficient. This is the most consequential
limitation of the current design.

\textbf{Structural fixity vs.\ adaptability.}
BN structure is fixed at deployment. Adding a new
evidence source requires redeployment and resets all
claim history. In a live healthcare setting where
evidence types evolve, this rigidity is a practical
barrier.

\textbf{Single claim variable encoding.}
The current design uses two independent root nodes
(\texttt{PPH} and \texttt{PPR}). A cleaner
alternative is a single binary root variable
$C\in\{0,1\}$, which would halve the number of root
configurations, reduce CPT size by half, and
eliminate the marginal independence assumption between
claim variables. We retain the two-root design for
compatibility with~\cite{bn-oracle} but recommend the
single-variable encoding for future deployments.

\subsection{Limitations}
\label{sec:limitations}

\textbf{Synthetic CPT parameters.}
All experiments use either neutral or heuristically
specified CPT values. Gas and fidelity results are
independent of CPT values by construction, but
posterior quality depends critically on parameter
correctness. Clinically grounded CPT elicitation from
real physician-visit data is necessary before any
clinical deployment.

\textbf{Limited BN size range.}
BN sizes beyond 32~evidence nodes are not tested. The
$O(N_e)$ initialization bound is consistent with
measurements at 4, 8, and 32~nodes but is not
validated beyond this range.

\textbf{Single oracle operator.}
The current prototype uses a single trusted oracle
operator. Contention, failure, and adversarial
behaviour of the oracle are not evaluated. Gas and
latency overhead of a multi-operator committee design
remain unmeasured.

\textbf{Privacy exposure.}
Committing binary evidence assignments to a public
ledger may expose patient-linked information. Under
GDPR and HIPAA, this requires either a permissioned
deployment or evidence commitment via cryptographic
commitment schemes before any live deployment.


%% ============================================================
\section{Conclusion}
\label{sec:conclusion}
%% ============================================================

This paper addresses a fundamental tension in
BN-based oracle design: on-chain inference preserves
verifiability but does not scale; off-chain inference
scales but severs the verifiability link.
LUNA resolves this by treating the blockchain as an
authenticated parameter and audit plane rather than
an inference engine. BN structure and CPTs are
anchored on-chain; inference executes off-chain; and
every evidence assignment, posterior output, and state
transition is committed as an immutable,
reconstructible transcript. Correctness is auditable
post-hoc from ledger state alone.

Four concrete results are established. First, per-query
on-chain gas cost is independent of BN size by
construction, confirmed empirically across 4 to
32~evidence nodes on Sepolia, with a
\placeholder{G}\%~reduction at 32~nodes relative to
full on-chain BN inference. Second, fixed-point
encoding under $\texttt{SCALE}=10^6$ introduces a
maximum absolute error bounded by $5\times10^{-7}$
per posterior, with no decision inconsistencies across
all tested configurations. Third, the claim lifecycle
supports incremental, out-of-order evidence arrival
across multiple concurrent claims with stable
per-update gas cost. Fourth, the full inference
transcript for any claim is reconstructible and
independently verifiable from on-chain state alone.

These results hold under explicit scope conditions:
BN structure is fixed at deployment, evidence
variables are binary, CPT parameters are synthetic
in this evaluation, and the oracle operator is a
single trusted party. These conditions define the
boundary of this contribution and the frontier for
future work.

Three directions are prioritised. The most
consequential is oracle centralization: a single
operator can submit incorrect posteriors detectable
only post-hoc. The immediate next step is a
multi-operator committee design with on-chain
detection of divergent submissions. The second
direction is clinical CPT validation: evaluating LUNA
with CPTs elicited from real physician-visit data is
necessary before any clinical deployment. The third
direction is privacy-preserving inference: integrating
zero-knowledge proofs for posterior verification
without evidence disclosure would simultaneously
address the privacy constraint and the runtime
correctness gap.

LUNA demonstrates that interpretable, auditable
BN-based oracle decisions are practically deployable
without gas costs that scale with model complexity.
The architecture is explicit about what it does not
yet provide---runtime correctness enforcement,
clinical parameter validation, and privacy
preservation---and these define the starting point
for the next stage of this work.


%% ============================================================
%% REFERENCES
%% ============================================================
\bibliographystyle{IEEEtran}
\bibliography{references}

\end{document}
